\subsection{Ambiente de Testes}

A validação do HormuzNet foi conduzida utilizando o \textit{Docker Compose} para simular um ambiente distribuído em rede local. Foram instanciados 4 \textit{brokers}, 6 drones autônomos e 20 sensores. O monitoramento foi realizado através do dashboard tático acessível via navegador.

\subsection{Análise de Resiliência e Consistência}

A Tabela~\ref{tab:resiliencia} apresenta os resultados dos testes de estresse e falhas injetadas no sistema para verificar a capacidade de recuperação da malha.

\begin{table}[ht]
  \centering
  \caption{Testes de resiliência e comportamento da malha.}
  \label{tab:resiliencia}
  \begin{tabular}{p{1.5cm}p{4.5cm}p{5.5cm}c}
    \toprule
    \textbf{ID} & \textbf{Cenário de Falha} & \textbf{Resposta do Sistema} & \textbf{Status} \\
    \midrule
    TR-01 & Queda do Broker NW & Drones NW realizam \textit{fallback} para Broker NE & OK \\
    TR-02 & Sensor Multicast & Alerta recebido simultaneamente por todos os brokers & OK \\
    TR-03 & Drone Abatido & Broker detecta perda e reenvia alerta para a fila global & OK \\
    TR-04 & Conflito Lamport & Brokers ordenam ocorrências identicamente via timestamps & OK \\
    TR-05 & Partição de Rede & Brokers isolados operam regionalmente e sincronizam ao retornar & OK \\
    \bottomrule
  \end{tabular}
\end{table}

\subsection{Desempenho da Sincronização}

As métricas de latência e processamento foram coletadas durante uma simulação de 10 minutos com alta taxa de ocorrências (1 evento a cada 2 segundos). Os resultados estão resumidos na Tabela~\ref{tab:performance}.

\begin{table}[ht]
  \centering
  \caption{Métricas de desempenho e sincronização.}
  \label{tab:performance}
  \begin{tabular}{p{5cm}p{3.5cm}p{4.5cm}}
    \toprule
    \textbf{Métrica} & \textbf{Valor Observado} & \textbf{Observação} \\
    \midrule
    Tempo de Despacho (médio) & $\approx$\,1.5 s & Da ocorrência ao envio do drone \\
    Divergência de Lamport & 0 ms & Filas mantiveram-se idênticas \\
    Uso de CPU por Broker & $<$ 5\% & Eficiência das goroutines \\
    Taxa de Sucesso de Missões & $\approx$\,91\% & Impacto do modelo de 10\% de falha \\
    Latência P2P entre Brokers & $<$ 10 ms & Comunicação TCP em rede local \\
    \bottomrule
  \end{tabular}
\end{table}

\subsection{Discussão dos Resultados}

Os experimentos demonstraram que a descentralização atingiu o objetivo de eliminar SPOFs. A queda de um \textit{broker} não interrompeu o monitoramento de sua região, pois o \textit{fallback} dos drones para nós vizinhos ocorreu em menos de 1 segundo. A consistência garantida pelo algoritmo de Lamport foi fundamental para evitar que múltiplos drones atendessem a mesma ocorrência, validando a eficácia da exclusão mútua distribuída implementada.

A principal limitação observada foi a escalabilidade do \textit{multicast} em redes que não suportam IGMP de forma eficiente, o que poderia restringir o sistema em ambientes de rede WAN complexos. Contudo, no escopo de uma rede de monitoramento regional, a solução mostrou-se extremamente robusta.
